\chapter{参考文献}

% ===== PDF p.102 =====
\begin{references}
\refitem{Art94}{Art94} James Arthur. \emph{Unipotent automorphic representations: conjectures}. Department of Mathematics, University of Toronto, 1994.
\refitem{Art02}{Art02} James Arthur. A note on the automorphic Langlands group. \emph{Canadian Mathematical Bulletin}, 45(4):466--482, 2002.
\refitem{Art05}{Art05} James Arthur. An introduction to the trace formula. \emph{Harmonic analysis, the trace formula, and Shimura varieties}, 4:1--263, 2005.
\refitem{BC09}{BC09} Olivier Brinon and Brian Conrad. \emph{CMI summer school notes on $p$-adic Hodge theory}, 2009.
\refitem{BG10}{BG10} Kevin Buzzard and Toby Gee. The conjectural connections between automorphic representations and Galois representations. \emph{Automorphic forms and Galois representations}, 1:135--187, 2010.
\refitem{Bor79}{Bor79} Armand Borel. Automorphic $L$-functions. In \emph{Automorphic forms, representations and $L$-functions} (Proc. Sympos. Pure Math., Oregon State Univ., Corvallis, Ore., 1977), Part 2, pages 27--61, 1979.
\refitem{Bor80}{Bor80} A. Borel. Stable and $L^2$-cohomology of arithmetic groups. \emph{Bull. Amer. Math. Soc. (N.S.)}, 3(3):1025--1027, 1980.
\refitem{BRa}{BRa} J. N. Bernstein and K. Rumelhart. Lectures on $p$-adic groups. 未出版, 见 \url{http://www.math.tau.ac.il/~bernstei/Publication_list}.
\refitem{BRb}{BRb} Don Blasius and J. Rogawski. Zeta functions of Shimura varieties. \emph{Motives (Seattle, WA, 1991)}, 525--571. In \emph{Proc. Sympos. Pure Math.}, volume 55.
\refitem{BT13}{BT13} Raoul Bott and Loring W. Tu. \emph{Differential forms in algebraic topology}, volume 82. Springer Science \& Business Media, 2013.
\refitem{Cas73}{Cas73} William Casselman. On some results of Atkin and Lehner. \emph{Mathematische Annalen}, 201(4):301--314, 1973.
\refitem{CGP15}{CGP15} Brian Conrad, Ofer Gabber, and Gopal Prasad. \emph{Pseudo-reductive groups}, volume 26 of New Mathematical Monographs. Cambridge University Press, Cambridge, second edition, 2015.
\refitem{Col11}{Col11} P. Colmez. \emph{Éléments d'analyse et d'algèbre (et de théorie des nombres)}, second ed. Éditions de l'École Polytechnique, Palaiseau, 2011.
\refitem{CS17}{CS17} Ana Caraiani and Peter Scholze. On the generic part of the cohomology of compact unitary Shimura varieties. \emph{Annals of Mathematics}, pages 649--766, 2017.
\refitem{DE14}{DE14} Anton Deitmar and Siegfried Echterhoff. \emph{Principles of harmonic analysis}. Springer, 2014.
\refitem{Del71a}{Del71a} P. Deligne. Travaux de Shimura, Sém. Bourbaki 389. \emph{Lecture Notes in Math}, 244, 1971.
\refitem{Del71b}{Del71b} Pierre Deligne. Formes modulaires et représentations $\ell$-adiques. In \emph{Séminaire Bourbaki vol. 1968/69, Exposés 347--363}, pages 139--172. Springer, 1971.
\refitem{Del73}{Del73} Pierre Deligne. Formes modulaires et représentations de $GL(2)$. In \emph{Modular functions of one variable II}, pages 55--105. Springer, 1973.
\refitem{Del75}{Del75} Pierre Deligne. Courbes elliptiques: formulaire d'après J. Tate. In \emph{Modular functions of one variable, IV} (Proc. Internat. Summer School, Univ. Antwerp, Antwerp, 1972), volume 476, pages 53--73, 1975.
\refitem{Del79}{Del79} Pierre Deligne. Variétés de Shimura: interprétation modulaire, et techniques de construction de modèles canoniques. In \emph{Automorphic forms, representations and $L$-functions} (Proc. Sympos. Pure Math., Oregon State Univ., Corvallis, Ore., 1977), Part 2, pages 247--289, 1979.
\refitem{DR73}{DR73} Pierre Deligne and Michael Rapoport. Les schémas de modules de courbes elliptiques. In \emph{Modular functions of one variable II}, pages 143--316. Springer, 1973.
\refitem{DS74}{DS74} Pierre Deligne and Jean-Pierre Serre. Formes modulaires de poids 1. In \emph{Annales scientifiques de l'École Normale Supérieure}, volume 7, pages 507--530, 1974.
\refitem{DS05}{DS05} Fred Diamond and Jerry Michael Shurman. \emph{A first course in modular forms}, volume 228. Springer, 2005.
\refitem{Fla79}{Fla79} Dan Flath. Decomposition of representations into tensor products. In \emph{Automorphic forms, representations and $L$-functions} (Proc. Sympos. Pure Math., Oregon State Univ., Corvallis, Ore., 1977), Part 1, pages 179--183, 1979.

% ===== PDF p.103 =====
\refitem{FM95}{FM95} Jean-Marc Fontaine and Barry Mazur. Geometric Galois representations. \emph{Elliptic curves, modular forms, \& Fermat's last theorem (Hong Kong, 1993), Ser. Number Theory, I}, pages 41--78, 1995.
\refitem{Fon94}{Fon94} Jean-Marc Fontaine. Le corps des périodes $p$-adiques. \emph{Astérisque}, 223:59--111, 1994.
\refitem{Gel75}{Gel75} Stephen S. Gelbart. \emph{Automorphic forms on adele groups}, number 83. Princeton University Press, 1975.
\refitem{GH19}{GH19} Jayce R. Getz and Heekyoung Hahn. \emph{An introduction to automorphic representations}. 见网页 \url{https://www.math.duke.edu/~hahn/GTM.pdf}, 2019.
\refitem{GS17}{GS17} Philippe Gille and Tamás Szamuely. \emph{Central simple algebras and Galois cohomology}, volume 165. Cambridge University Press, 2017.
\refitem{Har88}{Har88} Michael Harris. Automorphic forms and the cohomology of vector bundles on Shimura varieties. In \emph{Proceedings of the Conference on Automorphic Forms, Shimura Varieties, and $L$-functions}, Ann Arbor, pages 41--91, 1988.
\refitem{Lan17}{Lan17} Kai-Wen Lan. \emph{An example-based introduction to Shimura varieties}. Preprint, 2017.
\refitem{LR87}{LR87} R. P. Langlands and M. Rapoport. Shimuravarietäten und Gerben. \emph{J. Reine Angew. Math.}, 378:113--220, 1987.
\refitem{Lur}{Lur} Jacob Lurie. \emph{Tamagawa numbers via nonabelian Poincaré duality}. 讲义见 \url{http://www.math.harvard.edu/~lurie/282y.html}.
\refitem{Mil05}{Mil05} James S. Milne. Introduction to Shimura varieties. \emph{Harmonic analysis, the trace formula, and Shimura varieties}, 4:265--378, 2005.
\refitem{MW89}{MW89} Colette Mœglin and J.-L. Waldspurger. Le spectre résiduel de $GL(n)$. In \emph{Annales scientifiques de l'École normale supérieure}, volume 22, pages 605--674, 1989.
\refitem{Pan19}{Pan19} Lue Pan. The Fontaine--Mazur conjecture in the residually reducible case. \emph{arXiv preprint} arXiv:1901.07166, 2019.
\refitem{PS16}{PS16} Vincent Pilloni and Benoît Stroh. Arithmétique $p$-adique des formes de Hilbert: surconvergence, ramification et modularité. \emph{Astérisque}, 382:195--266, 2016.
\refitem{Rib77}{Rib77} Kenneth A. Ribet. Galois representations attached to eigenforms with nebentypus. In \emph{Modular functions of one variable V}, pages 18--52. Springer, 1977.
\refitem{Rib90}{Rib90} Kenneth A. Ribet. From the Taniyama--Shimura conjecture to Fermat's last theorem. \emph{Annales de la Faculté des sciences de Toulouse: Mathématiques}, Ser. 5, 11(1):116--139, 1990.
\refitem{Ser97}{Ser97} Jean-Pierre Serre. \emph{Abelian $\ell$-adic representations and elliptic curves}. CRC Press, 1997.
\refitem{Tho}{Tho} Jack Thorne. \emph{Topics in automorphic forms} (notes taken by Chao Li). 课程讲义见 \url{http://www.math.columbia.edu/~chaoli/docs/AutomorphicForm.html}.
\refitem{Thu82}{Thu82} William P. Thurston. Three dimensional manifolds, Kleinian groups and hyperbolic geometry. \emph{Bulletin of the American Mathematical Society}, 6(3):357--381, 1982.
\refitem{Wil95}{Wil95} Andrew Wiles. Modular elliptic curves and Fermat's last theorem. \emph{Annals of Mathematics}, 141(3):443--551, 1995.
\end{references}

\vspace{2em}
\noindent\textbf{作者地址}：Mathematisches Institut, Universität Bonn, Endenicher Allee 60, 53115 Bonn, Deutschland。电子邮件：\texttt{ja@math.uni-bonn.de}。
